\section{Actual radicals and a conditional ABC obstruction}
\label{sec:rt-integer-radical-gate}
The preceding sufficient conditions involve actual radical arithmetic. The
following elementary bridge is also checked against the standard logarithmic
$abc$ definition in the repository. Its family-existence premise remains open.

For coprime positive integers $a,b$, put $c=a+b$ and
\[
 M=a^2+ab+b^2,\qquad N=(ab)^2+abM+M^2.
\]
The two transformed triples are $(ab,M,c^2)$ and $(abM,N,c^4)$.
The next identities charge the new prime supports exactly; no upper bound
for them is assumed in the identities.

\begin{proposition}[Exact two-step radical transport]
\label{prop:rt-two-step-radical}
Both displayed triples are primitive, and
\begin{align*}
 \operatorname{rad}(abMc^2)
   &=\operatorname{rad}(abc)\operatorname{rad}(M),\\
 \operatorname{rad}(abMNc^4)
   &=\operatorname{rad}(abc)\operatorname{rad}(M)\operatorname{rad}(N).
\end{align*}
For any positive integers $d,V,Q,g$, with no disjoint-support requirement,
\[
 \operatorname{rad}(dVQ^g)\le dVQ.
\]
\end{proposition}
\begin{proof}
Modulo $a$ and $c$, the norm $M$ reduces to $b^2$, and modulo $b$ it
reduces to $a^2$. Hence $\gcd(M,abc)=1$. The identity $ab+M=c^2$
shows that the first transform is primitive. Apply the same argument to
the coprime pair $ab,M$; its norm $N$ is coprime to $abMc^2$ and therefore
to $abcM$. Radical multiplicativity for coprime factors and invariance
under positive powers give both equalities. For the last inequality use
$\operatorname{rad}(xy)\mid\operatorname{rad}(x)\operatorname{rad}(y)$,
$\operatorname{rad}(Q^g)=\operatorname{rad}(Q)$, and
$\operatorname{rad}(x)\le x$ for positive $x$.
\end{proof}

\begin{theorem}[An integer sufficient condition for a disproof]
\label{thm:rt-integer-gate}
If coprime positive seeds $(a,b)$ satisfying
\begin{equation}
 \bigl(\operatorname{rad}(M)\operatorname{rad}(N)\bigr)^5\le c
 \label{eq:rt-integer-gate}
\end{equation}
have unbounded $\log c$, then standard $abc$ is false.
\end{theorem}
\begin{proof}
Write $R_2=\operatorname{rad}(abMNc^4)$. Since
$\operatorname{rad}(abc)\le abc\le c^3$, the preceding proposition gives
\[
 R_2^5\le c^{15}
   \bigl(\operatorname{rad}(M)\operatorname{rad}(N)\bigr)^5
 \le c^{16}.
\]
Standard $abc$ with $\varepsilon=1/8$, applied to the actual primitive
triple $(abM,N,c^4)$, would provide one constant $C$ such that
\[
 4\log c\le \frac98\log R_2+C
 \le\frac{18}{5}\log c+C.
\]
Thus $(2/5)\log c\le C$, contradicting the stated unboundedness.
\end{proof}

\paragraph{Formal contract and open premise.}
The companion modules \texttt{RadicalTransport} and
\texttt{TwoStepABCObstruction} use Mathlib's radical and identify it with
the repository's product of distinct prime divisors. The last formal
statement has conclusion \texttt{not ABCConjecture}, conditional on the
explicit unbounded family in Theorem~\ref{thm:rt-integer-gate}. The source
definition of \texttt{ABCConjecture} is imported without alteration.
Neither the ordinary proof nor the formal implication constructs that
family. In particular, the implication is not a formal disproof of $abc$.
The new verification record states the checked source hashes, compiler,
Mathlib revision, and axiom audit separately from the earlier algebra-only
checkpoint.
